% compendium/execution-trading-costs/execution-scheduling.tex
% Entries migrated on 2026-10-15 from the legacy optimal-execution-problem blocks (slice 10), owned by
% the section overview execution-trading-costs; dispositions and source checks in notes/migration.md,
% the blocks verbatim in notes/artefacts/market-making-lob/derivations.tex.

\section{Execution scheduling and block orders}

\begin{qfmodel}[Mean--variance optimal execution: the Almgren--Chriss trade-off between impact and price risk]\label{mod:almgren-chriss-execution}
\qffield{Statement} An order is split over time to balance two forces: market impact pushes the schedule to trade slowly, price risk pushes it to trade fast \citep[app.~A, sec.~A.7, p.~274]{alpha-laruelle}. In Almgren and Chriss's discrete model a holding $X$ is liquidated over $N$ intervals of length $\tau=T/N$; $x_k$ is the holding at $t_k=k\tau$, with $x_0=X$ and $x_N=0$, and $n_k=x_{k-1}-x_k$ is sold in the $k$th interval \citep[author preprint of Dec.\ 2000, sec.~1.1, p.~7]{alpha-almgren}. The price follows $S_k=S_{k-1}+\sigma\tau^{1/2}\xi_k-\tau g(n_k/\tau)$, with permanent impact $g$; the $k$th sale executes at $\tilde S_k=S_{k-1}-h(n_k/\tau)$, with temporary impact $h$ that does not enter later prices (same preprint, secs.~1.2--1.3, eqs.~(1)--(2), pp.~8--9). The shortfall $XS_0-\sum_kn_k\tilde S_k$ has mean $E(x)$ and variance $V(x)=\sigma^2\sum_k\tau x_k^2$ (eqs.~(4)--(5), p.~10), and an efficient schedule minimises
\[ E(x)+\lambda V(x),\qquad \lambda\ \text{the risk aversion} \]
(sec.~2.1, eq.~(15), p.~13). With linear impact, $g(v)=\gamma v$ and $h=\epsilon\,\mathrm{sgn}(n_k)+(\eta/\tau)n_k$, the unique minimiser for $\lambda\ge0$ is
\[ x_j=\frac{\sinh\bigl(\kappa(T-t_j)\bigr)}{\sinh(\kappa T)}\,X,\qquad \frac2{\tau^2}\bigl(\cosh(\kappa\tau)-1\bigr)=\tilde\kappa^2=\frac{\lambda\sigma^2}{\tilde\eta},\quad\tilde\eta=\eta-\tfrac12\gamma\tau, \]
and $\theta=1/\kappa$ is the trade's half-life, independent of $T$ \citep[same preprint, secs.~2.2--2.3, eqs.~(16)--(17), pp.~13--15, page images of pp.~11 and~14]{alpha-almgren}.

\qffield{Notation} Lehalle and Laruelle's discrete version, used by the legacy notes, is for a buy order of $v^\star$ over $N$ slices: $v_n$ shares in slice $n$, market volume $V_n$, volatility $\sigma_n$ in currency per slice, and execution price $\tilde S_n(v_n)=S_n+\kappa\sigma_n(v_n/V_n)^{\gamma}$, with $\kappa$ and $\gamma$ stock-dependent constants \citep[app.~A, sec.~A.7.1, eqs.~(A.16)--(A.19), pp.~275--276]{alpha-laruelle}. Here $\kappa$ and $\gamma$ are impact parameters, not Almgren and Chriss's decay rate $\kappa$ and permanent-impact slope $\gamma$; $\lambda$ is risk aversion in both.

\qffield{Conditions} The price is a martingale over the horizon, with no view on its trend \citep[app.~A, sec.~A.7.2, p.~277]{alpha-laruelle}; volumes and volatilities per slice are taken as known in advance (p.~276); the schedule is static, which Almgren and Chriss show is optimal under their price dynamics (preprint, n.~10, p.~13); the trades do not change sign (p.~13). In Lehalle and Laruelle's version:
\begin{itemize}
\item \emph{Temporary impact only.} The concession $\kappa\sigma_n(v_n/V_n)^\gamma$ is in currency per share and applies only to slice $n$'s fills: an instantaneous execution cost, not a transient displacement of later prices. The slice costs $v_n$ times it; the expected impact cost of the order is $\kappa\sum_n\sigma_nv_n^{\gamma+1}/V_n^{\gamma}$, in currency, not per unit time (eq.~(A.20), p.~276). There is no permanent-impact term and no spread or fee term, unlike Almgren and Chriss's $g$ and $\epsilon$.
\item \emph{Expectation only} ($\lambda=0$). The schedule $v_n^{\exp}=v^\star\,(V_n/\sigma_n^{1/\gamma})\big/\sum_\ell V_\ell/\sigma_\ell^{1/\gamma}$ (Proposition~2, eq.~(A.22), p.~277), uniform if $V_n$ and $\sigma_n$ are constant (Proposition~3).
\item \emph{Variance only.} Trade everything in the first slice (Proposition~4, eq.~(A.25), p.~278); in Almgren and Chriss this gives zero variance ``because of the way that we have discretized time'' (preprint, eqs.~(12)--(13), p.~12).
\item \emph{Mean--variance}, $C=\mathbb E(W)+\lambda\operatorname{Var}(W)$ (eq.~(A.26), p.~279). For $\gamma=1$ the first-order conditions give a recurrence in the remaining quantity, solved by shooting on the first slice (eq.~(A.28), pp.~279--280); larger $\lambda$ trades faster, smaller $\lambda$ more linearly (Figure~A.11, pp.~280--281).
\end{itemize}

\qffield{Practice} A trading algorithm has two layers: a risk-control layer, at a time scale around 10 minutes, that bounds the trading rate from the schedule; and an opportunistic layer, such as smart order routers and liquidity seekers, that chooses the participation rate within those bounds and may react to alpha signals \citep[app.~A, sec.~A.7, p.~274]{alpha-laruelle}. The model's parameters are estimated daily or weekly (p.~274). Interpretation, checked numerically for this slice: the expectation-only schedule equalises marginal impact cost across slices; for $\gamma=1$ with constant $\sigma_n$ and $V_n$, the mean--variance schedule in the legacy notation equals Almgren and Chriss's $\sinh$ form with $\cosh k=1+\lambda\sigma V/(2\kappa)$ and Lehalle and Laruelle's recurrence; Lehalle and Laruelle's model is Almgren and Chriss's linear case without the schedule-independent terms (the permanent $\tfrac12\gamma X^2$ and the fixed cost $\epsilon X$ of a one-sided order) and without the $\tilde\eta$ correction, with $\kappa\sigma_n/V_n$ in place of $\eta/\tau$ and $\sigma_n^2$ in place of $\sigma^2\tau$.

\qffield{Pitfalls} Indexing. Lehalle and Laruelle's price step (A.16) prices slice $n$ after shock $n$, and their variance (A.24) uses the quantity remaining from slice $n$ on, yet they state the terminal conditions as $x_0=v^\star$, $x_N=0$; they also label Proposition~5 as ``variance minimization'' although it is the mean--variance result, and drop $\kappa$ from the intermediate Lagrange condition, which only rescales the multiplier. The legacy notes use the Almgren--Chriss timing instead: slice $n$ at $S_n$, then $S_{n+1}=S_n+\sigma_n\xi_{n+1}$, so the remaining quantity $x_n$ after slice $n$ bears the next shock and the risk term $\lambda\sum_nx_n^2\sigma_n^2$ is consistent (Interpretation). Comparative statics. The legacy objective said higher risk aversion front-loads the schedule ``while higher impact sensitivity back-loads it''. Interpretation: the minimiser depends on $\lambda/\kappa$ only, since $C-v^\star S_0=\kappa\bigl[\sum_n\sigma_nv_n^{\gamma+1}/V_n^\gamma+(\lambda/\kappa)\sum_nx_n^2\sigma_n^2\bigr]$; larger $\kappa$ moves the schedule towards the expectation-only one, not beyond it; in the cases checked the remaining quantity never exceeded the expectation-only path, which is not a proof. Shape. The legacy schedule taxonomy called the mean--variance path convex; that holds for constant parameters (the $\sinh$ form; Lehalle and Laruelle, p.~279), but with time-varying $V_n$ and $\sigma_n$ the path follows the volume profile and need not be convex (all 200 random cases checked were not). The parameters $\kappa,\gamma$ of a scheduling model are inputs; the empirical exponents of \hyperref[meth:almgren-impact-estimation]{Almgren et al.} (0.6, against size relative to volume during execution, with a 30-minute post-trade horizon) are estimates of a related but not identical quantity. The legacy taxonomy's ``dark pools and limits'' in the tactical overlay are not in the passages checked.

\qffield{See also} \hyperref[def:price-impact]{price impact} (temporary against transient impact; the Almgren--Chriss impact functions); \hyperref[meth:almgren-impact-estimation]{direct estimation of impact}; \hyperref[ev:intraday-impact-profile]{intraday impact profile} (participation and duration); \hyperref[res:no-price-manipulation]{absence of price manipulation} (conditions on impact models used in schedules); \hyperref[def:block-trade]{block trades}. The pathway \texttt{price-impact-execution} develops impact models and optimal schedules beyond this entry.

\qffield{Further reading} The legacy blocks it replaces, verbatim: \nolinkurl{notes/artefacts/market-making-lob/derivations.tex}. The journal version of Almgren and Chriss (2001) was not compared with the preprint.

\qffield{Index} optimal execution; Almgren--Chriss; mean--variance; implementation shortfall; risk aversion; execution schedule; trade half-life; expectation-minimising schedule; VWAP-like schedule; temporary impact cost; participation; trading algorithm layers

\qffield{Status} Verification: partly checked 2026-10-15 (against the text layer of a local copy of \texttt{alpha-laruelle}, 2018 World Scientific printing, app.~A, sec.~A.7, pp.~273--281; and \texttt{alpha-almgren} in the authors' preprint of December 2000, pp.~7--15, page images of pp.~9, 11 and~14; formulas and comparative statics checked numerically). Stability: draft. Currency: timeless (model).
\end{qfmodel}

\begin{qfdefinition}[Block trade: an order too large to fill by normal procedures, and the four problems of arranging it]\label{def:block-trade}
\qffield{Statement} A block trade results from an order too large to fill easily through normal trading procedures, typically more than a day's normal volume; most block traders think of a block as more than a quarter of a day's average volume in an actively traded stock, and in thinly traded instruments a few thousand shares can be one \citep[chap.~``Block Traders'', sec.~15.1, p.~322]{alpha-harris}. Block initiators face four problems \citep[same chap., sec.~15.2, pp.~323--327]{alpha-harris}:
\begin{itemize}
\item \emph{latent demand}: willing counterparties have not issued orders, because orders are costly to manage or because they have not yet thought about trading;
\item \emph{order exposure}: shopping the block lets front runners, same-side traders and opposite-side traders act on it, accelerating its price impact (Table~15-1, p.~325);
\item \emph{price discrimination}: liquidity suppliers fear that more size will follow, so they want the true size before offering liquidity;
\item \emph{asymmetric information}: suppliers suspect large, especially anonymous, traders of being informed, and demand concessions as in the adverse-selection spread.
\end{itemize}

\qffield{Conditions} Exchanges classify blocks by fixed sizes for statistics, the NYSE at 10{,}000 shares or more regardless of activity or price; Harris suggests classifying by a fraction of average daily volume instead \citep[same chap., sec.~15.1 and box, p.~323]{alpha-harris}. Interpretation: the legacy suggestion to measure a block by typical depth or participation capacity follows the same logic but is not in the passage.

\qffield{Practice} Solving the problems requires credible information about the initiator: that it wants to trade, exposure only to trustworthy traders, its full size, and evidence that it is uninformed. Since exchanges carry only price and size, most blocks are arranged by telephone upstairs through block brokers and dealers with reputations to protect \citep[same chap., secs.~15.2.2--15.2.5, pp.~324--327]{alpha-harris}. Crossing networks, hidden orders and order splitting are the alternatives when orders cannot be exposed (p.~325). \emph{Sunshine trading} (slice~34) is the opposite: announcing identity, intentions, full size and reasons for trading. It works when the sunshine trader is well known and known to be uninformed and honest, so in large markets at best for the largest traders, and fails when others suspect information or hidden size (``wolves in sheep's clothing''), because such motives cannot be verified on an exchange floor or screen \citep[same chap., sec.~15.3, pp.~327--328]{alpha-harris}; by revealing their trades, sunshine traders also give free trading options to order anticipators (p.~328). Block dealers and brokers audit their clients' motives: uninformed traders have utilitarian reasons unrelated to fundamental values, and block traders are most confident when the client has an easily verified reason to trade that instrument at that time; discretion over which instrument or when suggests strategic choice (Table~15-2: a firm liquidating a diversified portfolio to fund a project, probably uninformed; a legatee buying a single stock, probably informed) \citep[same chap., sec.~15.4.4, pp.~330--331]{alpha-harris}. Full and timely trade reporting favours block brokers over block dealers: a dealer who must report quickly can be front-run while liquidating the position, and reports help others infer whether the initiator was informed, which raises the importance of the motive audit; the weakest reporting is in dealer markets (sec.~15.4.3.1, p.~330). U.S. equity trades were reported within 90 seconds, and the London Stock Exchange had moved from seven days for the largest blocks to 90 minutes (box, p.~331, as of 2002). Harris also describes how blocks met the exchange book in U.S. listed equities in 2002 \citep[same chap., sec.~15.6.1, pp.~333--334]{alpha-harris}:
\begin{itemize}
\item brokers and dealers who were members of an exchange had to print their block trades in exchange-listed securities at an exchange if the trade took place while the market was open, presenting the matched orders so that the exchange's precedence rules protected all traders;
\item a \emph{size precedence} rule let matched crosses of more than 25{,}000 shares outsize the book at the trade price, jumping orders with time precedence, while orders standing at better prices were filled at the block price, since price priority was enforced universally;
\item Harris calls this a compromise between block traders, who want their crosses unbroken, and standing limit-order traders, who want a return on the liquidity they offer;
\item block traders often printed crosses at regional exchanges, whose books were thinner.
\end{itemize}

\qffield{Pitfalls} The legacy summary, ``disclose enough to find size, but not enough to be exploited'', is a fair reading of the tension between the latent-demand and order-exposure problems (Interpretation), but Harris's own resolution adds disclosure of identity and true size to trusted counterparties. The legacy price-discrimination item (``suppliers fear selling too cheaply or buying too richly'') paraphrases the fear that more size will follow. The legacy motive-audit process (slice 13) lists funding, diversification, hedging, benchmark change and speculation, and ``comparing the initiator's past behaviour''; Harris's examples are funding, diversification, covering a short, hedging and currency needs, and ``benchmark change'' and past behaviour are not in the passage. The legacy reporting caveat says that delayed reporting gives informed block traders room to exploit stale prices; Harris's argument is that timely reporting hurts facilitating dealers, and elsewhere that delayed reporting helps traders building or unwinding positions at the expense of other traders (chap.~``Buy-Side Traders'', p.~389). Some motive audits may be unreliable where prices are too volatile for block traders to build reputations for them (p.~331). The legacy market rule \emph{Equity Block Prints and Size Priority} (slice 18) merged three things that should be kept apart: where a trade is \emph{reported} or printed; its \emph{execution priority} against orders in the book (size against time precedence); and \emph{trade-through protection} of better-priced quotes on other venues, which in the United States is now Rule 611 (\hyperref[def:best-execution]{best execution}). Its ``trade-reporting facility'' and its size-priority exceptions are stated as current practice; only Harris's 2002 exchange rule was checked, and current venue rules are a backlog item of \texttt{market-structure-execution-quality}. Its design trade-off (protect the book too much and blocks leave; protect blocks too much and displayed liquidity loses value) restates Harris's compromise.

\qffield{See also} \hyperref[def:order-anticipation]{order anticipation} (front running exposed orders); \hyperref[def:order-exposure-decisions]{order exposure decisions}; \hyperref[def:dark-pools-crossing]{dark pools and crossing networks}; \hyperref[res:adverse-selection-spread-component]{spread components} (the adverse-selection concession); \hyperref[mod:almgren-chriss-execution]{mean--variance execution} (scheduling an order through the open market instead).

\qffield{Further reading} The legacy blocks it replaces, verbatim: \nolinkurl{notes/artefacts/market-making-lob/derivations.tex}.

\qffield{Index} block trade; block order; upstairs market; latent demand; order exposure; shopping the block; price discrimination; asymmetric information; block brokers; block dealers; block positioners; trading motive audit; utilitarian traders; trade reporting; delayed reporting

\qffield{Status} Verification: partly checked 2026-10-15 (against the text layer of a local copy of \texttt{alpha-harris}, 2002 Oxford printing, pp.~322--327, and pp.~328--331 and~389 added in slice 13, and pp.~333--334 added in slice 18; sunshine trading, pp.~327--328, added in slice 34). Stability: draft. Currency: timeless (concepts); the NYSE threshold, venues and reporting deadlines as of 2002.
\end{qfdefinition}
